\label{ch:introduction}

\section{Purpose and how to read this document}

This document is the second inference-side companion to the \emph{Theoretical
Foundations of \texttt{tpeanuts}} document: where that document develops the
generic neutrino-propagation physics implemented in \texttt{core/} and
\texttt{medium/}, and \emph{Inference at Daya Bay} documents a reactor
antineutrino application, this one documents a real solar-neutrino
application built on a different statistical foundation than either --
fitting the solar mixing angle $\theta_{12}$, the solar mass-squared
splitting $\Delta m^2_{21}$, and the total active $^8$B solar-neutrino flux
$\Phi_{8B}$ to the real public data of the Sudbury Neutrino Observatory's
salt phase (Phase~II). Every number quoted here that is not explicitly
attributed to a fitted \texttt{tpeanuts} result is a real, published,
externally verifiable quantity, traceable through the References chapter.

The five chapters are organised as follows: this chapter introduces the
experiment itself -- the salt-phase detector modification, its three
detection reactions, and the data-taking timeline -- with no reference yet to
\texttt{tpeanuts}; \cref{ch:data-release} catalogues the real official
published tables this document's data was transcribed from, table by table,
including an honest account of which entries could and could not be read with
full confidence; \cref{ch:detector-implementation} is the physics-and-mathematics
core, developing the equivalent-flux forward model and the day/night Earth
regeneration \texttt{tpeanuts.detector.sno\_ii} implements;
\cref{ch:inference-process} documents the correlated-covariance statistical
machinery -- including a genuine extension to this project's own shared
inference code, built specifically to support this document -- and the
fixed-slice/profiled likelihood-scan distinction; \cref{ch:results} presents,
interprets, and summarises the numerical results obtained from
\texttt{inference6\_sno\_ii.ipynb}, compared directly against SNO's own
published best fit and likelihood contours.

\section{Overview}

The Sudbury Neutrino Observatory (SNO) is a heavy-water Cherenkov detector
located 2039~m underground in the Creighton mine near Sudbury, Ontario,
Canada, built to resolve the decades-old solar neutrino problem by measuring,
independently and simultaneously, both the $\nu_e$-only flux and the
total-active-flavour flux of $^8$B solar neutrinos. Its original (``Phase~I''
or ``pure D$_2$O'') running period established, via a pure-D$_2$O
neutral-current measurement, that solar $\nu_e$ genuinely change flavour en
route from the Sun rather than simply disappearing
\parencite{Ahmad2002NCDiscovery}. This document concerns SNO's second running
phase, the \emph{salt phase} (Phase~II): between June 2001 and October 2003,
approximately 2000~kg of NaCl were dissolved into SNO's $\sim$1000 tonnes of
D$_2$O target, enhancing the detector's sensitivity to the neutral-current
reaction specifically \parencite{Aharmim2005SaltPhase}. The data set this
document's inference is built against is the full 391-live-day salt-phase
sample, with its real published day/night electron-recoil spectrum and
integrated fluxes \parencite{Aharmim2005SaltPhase}.

\begin{notebox}
\textbf{Scope of this document.} \texttt{tpeanuts} reproduces SNO's own real
\emph{equivalent-flux} observable definitions (Eq.~A1 of the primary source),
the real salt-phase energy-response parametrisation, real day/night
Earth-matter regeneration, and a real correlated statistical+systematic
covariance built from the collaboration's own published correlation and
sensitivity tables. It does \emph{not} reproduce SNO's own underlying
event-by-event maximum-likelihood signal extraction (the process that turned
raw PMT hit patterns into the 38 published observables in the first place --
see \cref{sec:salt-phase-principle}): that extraction is consumed as
already-published, external input, the same ``use the officially extracted
observable, not the raw data'' pattern this project already uses for
Daya Bay's published IBD spectra and IceCube's published detector-systematics
hypersurfaces. \Cref{ch:results} reports, honestly and without adjustment,
where two specific systematic-table entries could not be transcribed with
confidence from the published source and were conservatively set to zero
rather than guessed.
\end{notebox}

\section{The salt-phase detector modification}
\label{sec:salt-phase-detector}

\begin{theorybox}
SNO's heavy-water target is instrumented by $\sim$9500 photomultiplier tubes
mounted on a geodesic support structure surrounding a transparent acrylic
vessel, itself submerged in ordinary (light) water for shielding, following
the same general Cherenkov-detector design principle used by Super-Kamiokande
and (with GdLS instead of pure D$_2$O) Daya Bay. Dissolving NaCl into the
D$_2$O target changed the physics of neutron detection specifically, since
the dominant neutron-capture nucleus becomes $^{35}$Cl rather than deuterium
itself:
\begin{equation}
  n + {}^{35}\mathrm{Cl} \rightarrow {}^{36}\mathrm{Cl} + \gamma\text{'s}
  \quad(\text{total }\gamma\text{ energy}\approx 8.6~\mathrm{MeV}),
\end{equation}
compared to Phase~I's single $\approx6.25$~MeV $\gamma$ from
$n+d\rightarrow t+\gamma$. Three real, direct consequences follow
\parencite{Aharmim2005SaltPhase}: (i) the neutron-capture efficiency
increases by roughly a factor of three; (ii) the higher-multiplicity,
higher-energy $^{35}$Cl capture gamma cascade produces a measurably more
\emph{isotropic} PMT hit pattern than a single-gamma CC or ES Cherenkov cone,
providing a real, event-by-event discriminating observable
(\cref{sec:salt-phase-principle}); (iii) the NC capture-gamma energy spectrum
shifts to overlap the CC electron-recoil spectrum's own range more heavily
than in Phase~I, reducing (not increasing) how much spectral shape alone can
separate CC from NC -- the direct physical reason SNO's own salt-phase
analysis relies on event isotropy for that separation rather than
recoil-energy shape alone.
\end{theorybox}

\begin{notebox}
The detector depth (2039~m) used throughout \texttt{tpeanuts.detector.sno\_ii}
and \texttt{tpeanuts.detector.sno} alike is SNO's commonly quoted approximate
value, not independently re-verified here against a primary SNO
instrumentation paper -- its effect on Earth-crossing geometry is tiny
relative to the Earth's $\sim$6371~km radius, the same reasoning already
applied to Daya Bay's/Borexino's own approximate depths elsewhere in this
project.
\end{notebox}

\section{Three simultaneous detection reactions}
\label{sec:three-reactions}

\begin{theorybox}
SNO measures $^8$B solar neutrinos through three physically distinct
reactions on the same D$_2$O(+NaCl) target, occurring simultaneously in the
same data set:
\begin{align}
  \nu_e + d &\to p + p + e^- &&\text{(charged current, CC, } \nu_e\text{-only)}, \\
  \nu_x + e^- &\to \nu_x + e^- &&\text{(elastic scattering, ES, all flavours)}, \\
  \nu_x + d &\to p + n + \nu_x &&\text{(neutral current, NC, flavour-blind)}.
\end{align}
CC is sensitive only to the electron flavour, giving it a direct handle on
the survival probability $P_{ee}$; ES is sensitive to every active flavour
but with unequal couplings ($\sigma_e\approx6$--$7\times\sigma_{\mu,\tau}$ in
this energy range, since $\nu_e$ scatters via both charged- and
neutral-current diagrams while $\nu_\mu,\nu_\tau$ scatter via neutral current
only); NC has an equal cross section for every active flavour, so in the
Standard Model (three active flavours, unitary leptonic mixing) the NC rate
is exactly the total active $^8$B flux, independent of the oscillation
parameters -- SNO's own direct statement: ``the flux of all active neutrino
types that would produce the same NC rate''
\parencite{Aharmim2005SaltPhase}. This three-reaction combination is what let
SNO measure $\Phi_{8B}$ (via NC) and $P_{ee}(E)$ (via CC/NC and ES)
simultaneously, without relying on an external Standard Solar Model
flux prediction the way a single-reaction (CC-only) experiment must.
\end{theorybox}

\section{Statistical separation via event isotropy}
\label{sec:salt-phase-principle}

\begin{theorybox}
Unlike Phase~I -- where CC, ES, and NC contributions were separated primarily
by the shape of the reconstructed electron-energy spectrum, since the NC
capture-gamma peak sat at a distinctly lower energy than the bulk of the CC
spectrum -- SNO's own salt-phase analysis performs an event-by-event
statistical separation via an extended maximum-likelihood fit to four
reconstructed quantities per event: kinetic energy $T_{\rm eff}$, volume-weighted
radius $\rho$, an isotropy parameter $\beta_{14}$ (built from Legendre
polynomials of the angles between pairs of PMT hits, small for a
directional Cherenkov cone, larger for the more isotropic multi-gamma NC
cascade), and the angle to the Sun $\cos\theta_\odot$ (peaked forward for CC
and ES, flat for NC). This extraction -- performed entirely by the SNO
collaboration on their own raw PMT data, never by \texttt{tpeanuts} -- is
what produces the 38 published pseudo-observables this document's forward
model predicts and fits against (\cref{sec:38-observables}).
\end{theorybox}

\begin{notebox}
\texttt{tpeanuts} does not, and cannot, reconstruct this event-by-event
statistical separation: the underlying multi-dimensional probability density
functions ($P(T_{\rm eff},\beta_{14},\rho,\cos\theta_\odot)$ per reaction) were
never published. \texttt{detector.sno\_ii} instead predicts the \emph{result}
of that separation directly (\cref{ch:detector-implementation}), the same
principled choice already made for Daya Bay's IBD spectra and IceCube's
detector-systematics hypersurfaces elsewhere in this project.
\end{notebox}

\section{Data-taking period and physics goals}

\begin{theorybox}
The salt phase spans June 2001 to October 2003; this document's data set is
SNO's own full 391-live-day sample from that period, split into 128.5 days
of solar daytime and 262.9 days of solar nighttime livetime (their own
zenith-livetime table, \cref{ch:data-release}). The salt phase's own
headline results, quoted directly for later comparison
(\cref{ch:results}): total active $^8$B flux
$\Phi_{8B}=4.94^{+0.21}_{-0.21}({\rm stat})^{+0.38}_{-0.34}({\rm syst})\times10^6\,{\rm
cm^{-2}s^{-1}}$ (from the NC channel alone, independent of any oscillation
assumption); best-fit oscillation parameters, combined with the earlier pure
D$_2$O phase and (separately) with KamLAND reactor data,
$\Delta m^2=8.0^{+0.6}_{-0.4}\times10^{-5}\,{\rm eV}^2$,
$\theta=33.9^{+2.4}_{-2.2}$~degrees \parencite{Aharmim2005SaltPhase}. This
document's own fit (\cref{ch:results}) uses only the salt-phase data itself,
with no KamLAND or Phase~I combination -- a deliberately narrower, real
comparison baseline than SNO's own headline combined result.
\end{theorybox}
